

The groups $H^*(E_{12};\mathbb Z)$ are torsion free, with
the ranks in \eqref{maps:eq:E.S44}. For the actual fibre $Y$, the only
torsion in cohomology is
\[
       H^3(Y;\mathbb Z)_{\rm tors}
       =H^7(Y;\mathbb Z)_{\rm tors}=(\mathbb Z/6)^2.
 \tag{E.S48}\label{maps:eq:E.S48}
\]
Equivalently, its integral homology has the ranks \eqref{maps:eq:E.S33}
and torsion $(\mathbb Z/6)^2$ exactly in degrees two and six.
The twelfth return $T=\psi^{12}$ acts as the identity on
all integral cohomology except in degrees four and five.
In each of those two free groups it has, in an integral
basis, the form
\[
             \begin{pmatrix}1&72\\0&1\end{pmatrix}
                                  \oplus I_{691}.
 \tag{E.S49}\label{maps:eq:E.S49}
\]

For $n\geq1$, let $P_{12n}$ be the connected cyclic cover
of $P$ of degree $12n$ defined by $\chi$. Its Betti numbers,
independent of $n$, are
\[
       (1,13,133,528,1099,1384,1099,528,133,13,1).
 \tag{E.S50}\label{maps:eq:E.S50}
\]
Its integral homology torsion is exactly
\[
 H_j(P_{12n};\mathbb Z)_{\rm tors}=
 \begin{cases}
       (\mathbb Z/6)^2,&j=2,3,6,7,\\
       \mathbb Z/(72n),&j=4,5,\\
       0,&\text{otherwise}.
 \end{cases}
 \tag{E.S51}\label{maps:eq:E.S51}
\]
In particular these covers are pairwise nonhomotopy-equivalent.
These integral conclusions still do not recognize the smooth
$E_{12}$ or replace its boundary marking by homological data.

\begin{proof}
\emph{The integral cohomology of the actual cover.}
Apply the proof of \eqref{maps:eq:E.S42} with the coefficient lattice
$\mathbb Z[C_{12}]$, using its regular permutation matrix
$P_{12}$. In vertical degree $q$ set
\[
 \begin{aligned}
 T_{1,q}=P_{12}^4\otimes\bigwedge^q A_1^{-t},\quad
 &T_{2,q}=P_{12}^{-3}\otimes\bigwedge^q A_2^{-t},\\
 T_{0,q}=I_{12}\otimes\bigwedge^q M_0^{-t},
 &V_q=\mathbb Z^{12\binom4q}.
 \end{aligned}
 \tag{E.S52}\label{maps:eq:E.S52}
\]
This coefficient system computes $H^*(E_{12};\mathbb Z)$.
There is no integral averaging step. At the first finite
cap the cover has four components, each the ordinary
degree-three torus root cover; at the second there are
three ordinary degree-four torus root covers. An invariant
vector is determined by its value on one sheet of such
an orbit, with the other values its integral monodromy
translates. This identifies cap cohomology with the full
primitive invariant lattice. The cusp has twelve unchanged
components and the integral specialization of \eqref{maps:eq:B.5.16}.
Consequently its direct-image sheaf is exactly the integral
$j_*L_q$, including at every critical point.

The three graph-cocycle values again surject onto the
three local coinvariant groups, integrally. Thus $H^2(j_*L_q)=0$,
and its remaining cohomology is the kernel and cokernel
of the exact lattice map
\[
 \begin{gathered}
 V_q\longrightarrow
        \bigoplus_{i=0}^2\operatorname{im}(T_{i,q}-I),\\
           v\longmapsto((T_{i,q}-I)v)_i.
 \end{gathered}
 \tag{E.S53}\label{maps:eq:E.S53}
\]
Use Hermite column bases for the three image lattices;
this retains their integral indices rather than replacing
them by rational spans. The source and target ranks,
and the rank of \eqref{maps:eq:E.S53}, are respectively
\[
 \begin{array}{c|rrrrr}
 q&0&1&2&3&4\\ \hline
 \operatorname{rank}V_q&12&48&72&48&12\\
 \operatorname{rank}\bigoplus\operatorname{im}(T_{i,q}-I)
                      &17&92&126&92&17\\
 \operatorname{rank}(\text{E.S53})&11&47&71&47&11
 \end{array}
 \tag{E.S54}\label{maps:eq:E.S54}
\]
Every nonzero Smith entry is one. Hence $H^0(j_*L_q)$
has rank one and $H^1(j_*L_q)$ is free with ranks
$(6,45,55,45,6)$. The two-column Leray filtration has
free end groups and therefore splits as abelian groups.
This proves the torsion-free assertion and \eqref{maps:eq:E.S44}
integrally. The full integer matrices and the exact
Hermite/Smith check are in
\texttt{maps\_shifted\_cover\_integral.py}.

\emph{Why the marked rational images now determine the integral images.}
Each piece $A=E_{12}\times T^3$ in \eqref{maps:eq:E.S29} has free
integral homology. By Poincare--Lefschetz duality,
$H^{k+1}(A,\partial A;\mathbb Z)$ is free. The cokernel
of $H^k(A;\mathbb Z)\to H^k(\partial A;\mathbb Z)$
embeds in that free group, so the boundary image is a
primitive lattice. Its rational span has already been
computed in \eqref{maps:eq:E.S38}. Therefore the actual integral
image is precisely the saturation of that displayed
image in $\bigwedge^k\mathbb Z^8$.

Choose saturated column bases $\overline A_k,\overline B_k$
for the two images. The combined matrix
$[\overline A_k\ \overline B_k]$ has the ranks in the
second row of \eqref{maps:eq:E.S40}. All its nonzero Smith entries are
one except that, in each of degrees two and six, exactly
two entries are six. The Mayer--Vietoris exact sequence
has a free kernel term on its right, since all piece
cohomology is free. Thus its abelian extension splits,
and these Smith entries give exactly \eqref{maps:eq:E.S48}. The free
ranks remain \eqref{maps:eq:E.S33}.

\emph{The integral variation, including possible torsion components.}
Formula \eqref{maps:eq:E.S37} holds integrally. One can use the singular
prism of the loop of translations instead of differential
forms: the action map $S^1\times T^8\to T^8$ has
cohomology formula
\[
 a\longmapsto 1\otimes a+[S^1]^*\otimes\iota_v a.
\]
The difference between the two collar homotopies is
exactly this loop, so its connecting class is
$\partial\iota_v r(a)$. This proves the integral formula
also on torsion and still gives $(T^*-I)^2=0$.

For complete lattice control, take an integral kernel
basis of $[\overline A_k\ -\overline B_k]$. Its first
block, multiplied by $\overline A_k$, is a basis of the
common boundary lattice. Contract it by $v$ and reduce
modulo $[\overline A_{k-1}\ \overline B_{k-1}]$ using
the integral Smith transforms. In degrees four and five
the result has rank one and unique nonzero Smith entry
$72$. In degree three its two possible components
modulo six are both zero; in degree seven the common
boundary image is zero. All the remaining reductions
are zero. These computations are performed, with the
unimodular transforms verified, in
\texttt{maps\_shifted\_fibre\_integral.py}.
In degrees four and five the connecting cokernel is
free and embeds as a primitive summand of $H^k(Y)$,
because the following Mayer--Vietoris kernel is free.
Thus $72$ is the divisibility in the whole cohomology
group, not just in a chosen sublattice.

A rank-one integral map of Smith entry $72$ has the
form $72u\ell$, where $u$ and $\ell$ are primitive.
Its square is zero precisely when $\ell(u)=0$.
Extend $u$ to a basis of $\ker\ell$ and choose a vector
on which $\ell$ is one. This gives \eqref{maps:eq:E.S49}. The same
calculation proves identity on every other integral
cohomology group, including its torsion.

\emph{The finite-cover groups and the two potentially nontrivial extensions.}
The cover $P_{12n}$ is the mapping torus of $T^n$.
Since $(T^*-I)^2=0$, its variation is $n(T^*-I)$,
with Smith entry $72n$ in degrees four and five.
For every $k$ the integral Wang sequence gives
\[
 0\longrightarrow\operatorname{coker}(T^{n*}-I)_{k-1}
 \longrightarrow H^k(P_{12n};\mathbb Z)
 \longrightarrow\ker(T^{n*}-I)_k\longrightarrow0.
 \tag{E.S55}\label{maps:eq:E.S55}
\]
The ranks immediately give \eqref{maps:eq:E.S50}. Whenever the right
group is free, \eqref{maps:eq:E.S55} splits as an abelian sequence.
In particular degrees four and eight have torsion
$(\mathbb Z/6)^2$, contributed by $H^3(Y)$ and $H^7(Y)$
respectively, while degrees five and six have torsion
$\mathbb Z/(72n)$. Every other such degree has no torsion.

The only right groups in \eqref{maps:eq:E.S55} that are not free
occur in degrees three and seven. Their extensions
must not be split by assumption. Instead use the
nondegenerate integral torsion linking pairing of the
closed oriented ten-manifold $P_{12n}$: cohomological
torsion in degree $k$ is Pontryagin dual to that in
degree $11-k$. The already fixed torsion in degree
eight therefore forces the degree-three torsion to
be $(\mathbb Z/6)^2$, and the fixed degree-four torsion
forces the same answer in degree seven. The free
ranks are already known. This resolves precisely
the two possibly absorbing extensions. The universal
coefficient theorem now translates these cohomology
groups to \eqref{maps:eq:E.S51}. The order $72n$ in the middle
groups distinguishes the covers for distinct $n$.
\end{proof}


